% ================================================================
%  book-meta.tex — everything that names this book, in one place.
%
%  [2026-11-13] The cover says "Python from Zero". The on-ramps call it
%  "Programming Language From Zero: Python", and that difference is deliberate:
%  the topic is the language, Python is the instance. Do not make them match.
%  See CLAUDE.md, "The book's title, and the name the on-ramps give it".
% ================================================================
\newcommand{\BookTitle}{Python from Zero}
\newcommand{\BookTitleLines}{Python\\from Zero}
\newcommand{\BookSubtitle}{Write It, Understand It, Build With It}
\newcommand{\BookLongSubtitle}{From Your First Program to a Tested,\\
  Typed Project Somebody Else Would Accept}
\newcommand{\BookTagline}{From your first program to code you can defend}
\newcommand{\BookAuthor}{Lakshmi Deepak}
\newcommand{\BookCopyrightHolder}{Lakshmi Deepak}
\newcommand{\BookEdition}{First Edition}
\newcommand{\BookYear}{2026}
\newcommand{\BookRevisions}{2026-11-13 & Contents approved; writing begins \\}
\newcommand{\BookTestedWith}{Python 3.14 on macOS 26}
\newcommand{\BookCodeNote}{}
\newcommand{\BookCoverImage}{}
\newcommand{\BookCategory}{Programming / Python}
\newcommand{\BookPDFSubject}{Learning Python from the first program to a tested, typed,
  packaged project, and the data and concurrency a real program needs.}
\newcommand{\BookKeywords}{python, testing, type hints, packaging, concurrency, sql}
\newcommand{\BookBlurb}{%
  Most people learn enough Python to make something work, and then stop. This book takes you
  past that. It teaches the language twice on purpose: first so you can write it, then so you
  can say why it behaves as it does --- and then it puts that language to work against real
  data, real concurrency and real failure.

  \textbf{You'll learn how to:}
  \begin{itemize}
    \item Write Python to a standard somebody else would accept, and let a tool prove it
    \item Test what you wrote, type it, package it, and run it on a machine that is not yours
    \item Explain what Python is doing, from names and objects to decorators and the MRO
    \item Choose between threads, processes and async, and measure the choice
    \item Store data in a database, write the join, and find out why a query is slow
  \end{itemize}}
